% Definieren Sie hier Ihre Glossar-Einträge anhand der DHGE Beispiele.
% Diese Datei wird nur verwendet, wenn in der config.tex CGLOSSAR auf 1 steht
% (Voreinstellung). Bei CGLOSSAR=0 gilt stattdessen abk.tex (acro-Paket).
% Das Glossar ersetzt das Abkürzungsverzeichnis: Es enthält alle Fachbegriffe
% UND Abkürzungen, deren Bedeutung nicht für jeden Leser klar und eindeutig ist
% (auch firmeninterne Begriffe!).
% Das Glossar wird automatisch alphabetisch sortiert und formatiert
% (Fachbegriff fett auf eigener Zeile, Erläuterung eingerückt darunter).
%
% Fachbegriff:
%   \newglossaryentry{id}{
%       name = {Begriff},
%       description = {Erläuterung in einem (kurzen) Absatz},
%   }
%
% Abkürzung:
%   \newacronym[description = {Erläuterung}]{id}{ABK}{ausgeschrieben}
%
% Im Text verwenden Sie dann einfach \gls{id} (bzw. \ac{id}).
% LaTeX kümmert sich um den Rest.

\newglossaryentry{latex}{
    name = {LaTeX},
    description = {Ein Textsatzsystem zur Erstellung wissenschaftlicher Dokumente auf Basis des Satzprogramms \TeX. In diesem Template werden alle Formatierungen über LaTeX-Befehle vorgenommen.},
}

\newacronym[description = {Die Duale Hochschule Gera-Eisenach (DHGE) ist eine staatliche Hochschule des Freistaats Thüringen mit Standorten in Gera und Eisenach, an der duale Studiengänge in Kooperation mit Praxispartnern angeboten werden.}]
{dhge}{DHGE}{Duale Hochschule Gera-Eisenach}
